\documentclass[11pt]{article}
\usepackage[margin=1in]{geometry}
\usepackage{amsmath, amssymb, amsthm}
\usepackage{enumitem}
\usepackage{xcolor}
\usepackage{tcolorbox}
\tcbuselibrary{skins, breakable}

% Define color boxes for different framework sections
\newtcolorbox{problemconfigbox}{
    colback=blue!5, colframe=blue!40!black, title=PROBLEM CONFIGURATION ANALYSIS,
    fonttitle=\bfseries, breakable
}

\newtcolorbox{strategicbox}{
    colback=pink!5, colframe=pink!40!black, title=STRATEGIC FRAMEWORK APPLICATION,
    fonttitle=\bfseries, breakable
}

\newtcolorbox{tacticalbox}{
    colback=green!5, colframe=green!40!black, title=TACTICAL METHODOLOGY SELECTION,
    fonttitle=\bfseries, breakable
}

\newtcolorbox{conversionbox}{
    colback=yellow!5, colframe=yellow!40!black, title=CONVERSION TECHNIQUES APPLICATION,
    fonttitle=\bfseries, breakable
}

\newtcolorbox{verificationbox}{
    colback=red!5, colframe=red!40!black, title=VERIFICATION STRATEGY DESIGN,
    fonttitle=\bfseries, breakable
}

\newtcolorbox{roadmapbox}{
    colback=cyan!5, colframe=cyan!40!black, title=SOLUTION ROADMAP,
    fonttitle=\bfseries, breakable
}

\newtcolorbox{competitionbox}{
    colback=purple!5, colframe=purple!40!black, title=COMPETITION STRATEGY INTEGRATION,
    fonttitle=\bfseries, breakable
}

\newtcolorbox{keyinsightbox}{
    colback=orange!5, colframe=orange!40!black, title=KEY INSIGHT,
    fonttitle=\bfseries,breakable
}

\newtcolorbox{warningbox}{
    colback=red!10, colframe=red!60!black, title=WARNING,
    fonttitle=\bfseries,breakable
}

\newtcolorbox{tipbox}{
    colback=green!10, colframe=green!60!black, title=PRO TIP,
    fonttitle=\bfseries,breakable
}

\title{\textbf{2016 AMC 12B Problem 12: Strategic Analysis}}
\author{Using AMC12/COMC Problem-Solving Framework}
\date{}

\begin{document}
\maketitle

\section*{Problem Statement}
All the numbers $1, 2, 3, 4, 5, 6, 7, 8, 9$ are written in a $3\times3$ array of squares, one number in each square, in such a way that if two numbers are consecutive then they occupy squares that share an edge. The numbers in the four corners add up to $18$. What is the number in the center?

\textbf{Answer Choices:} 
$\textbf{(A)}\ 5\qquad\textbf{(B)}\ 6\qquad\textbf{(C)}\ 7\qquad\textbf{(D)}\ 8\qquad\textbf{(E)}\ 9$

\begin{problemconfigbox}
\subsection*{A. Problem Classification}
\begin{itemize}[leftmargin=*]
    \item \textbf{Problem Type}: To Find - determine the specific value in the center position
    \item \textbf{Difficulty Band}: Mid-band (Problem 12 of 25) - requires insight but not advanced techniques
    \item \textbf{Competition Context}: AMC12 (75 min, 25 problems) - approximately 3 minutes allocated
    \item \textbf{Mathematical Domain}: Combinatorics with parity/graph theory elements
\end{itemize}

\subsection*{B. Problem Structure Identification}
\begin{itemize}[leftmargin=*]
    \item \textbf{Given Information}: 
    \begin{itemize}
        \item Numbers $1-9$ arranged in $3\times3$ grid (one per square)
        \item Consecutive numbers must share an edge (not diagonal)
        \item Corner numbers sum to $18$
    \end{itemize}
    \item \textbf{Constraints}: 
    \begin{itemize}
        \item Consecutive number adjacency constraint (forms a path structure)
        \item Fixed sum constraint for corners
        \item All numbers $1-9$ used exactly once
    \end{itemize}
    \item \textbf{Objective}: Find the value at the center position
    \item \textbf{Hidden Assumptions}: 
    \begin{itemize}
        \item The arrangement must form a valid Hamiltonian path through the grid
        \item The center position has special properties (4 neighbors vs 2 for corners)
    \end{itemize}
    \item \textbf{Answer Choice Leverage}: Choices range from 5-9 and include both odd (5,7,9) and even (6,8) numbers. Parity analysis will help eliminate impossible options. The even choices provide a clue that checkerboard coloring may be relevant
\end{itemize}

\subsection*{C. Strategic Orientation}
\begin{itemize}[leftmargin=*]
    \item \textbf{Hypothesis}: The consecutive constraint creates a path structure that alternates positions
    \item \textbf{Conclusion}: Determine which of $\{5,6,7,8,9\}$ occupies the center
    \item \textbf{Notation}: Let the grid be:
    \[
    \begin{array}{|c|c|c|}
    \hline
    a & b & c \\
    \hline
    d & \boxed{x} & e \\
    \hline
    f & g & h \\
    \hline
    \end{array}
    \]
    where $x$ is the center value and $a+c+f+h=18$
    \item \textbf{Intuitive Approach}: 
    \begin{itemize}
        \item The consecutive constraint suggests analyzing parity and graph coloring
        \item Corner sum constraint provides arithmetic restriction
    \end{itemize}
    \item \textbf{Cross-Domain Potential}: 
    \begin{itemize}
        \item Graph Theory: Hamiltonian path on grid graph
        \item Parity Analysis: Checkerboard coloring argument
        \item Number Theory: Sum constraints and odd/even properties
    \end{itemize}
\end{itemize}
\end{problemconfigbox}

\begin{strategicbox}
\subsection*{A. Psychological Strategies}
\begin{itemize}[leftmargin=*]
    \item \textbf{Mental Toughness}: This problem requires recognizing the parity pattern - if first attempt doesn't work, try visualization with checkerboard coloring
    \item \textbf{Creativity Cultivation}: Think about what the consecutive constraint implies for position types (corner vs edge vs center)
    \item \textbf{Iterative Refinement}: Start with parity observation, then refine to specific arithmetic constraints
\end{itemize}

\subsection*{B. Investigation Strategies}
\begin{itemize}[leftmargin=*]
    \item \textbf{Startup Orientation}: Read carefully - "consecutive numbers share an edge" means a path structure
    \item \textbf{Get Your Hands Dirty}: Color the grid like a checkerboard and observe the alternating pattern
    \item \textbf{Penultimate Step}: Before finding center value, determine which positions (corners/edges/center) can hold odd vs even numbers
    \item \textbf{Wishful Thinking}: Wish that parity uniquely determines position types
    \item \textbf{Make It Easier}: First ignore the sum constraint and just analyze the consecutive constraint
    \item \textbf{Conjecture via Small Cases}: Try placing 1 in a corner and see what constraints emerge
    \item \textbf{Visualization}: Draw a checkerboard pattern on the grid
\end{itemize}

\subsection*{C. Late-Band Strategic Playbook}
\begin{itemize}[leftmargin=*]
    \item \textbf{Not applicable} - this is mid-band Problem 12, but techniques remain useful:
    \item \textbf{Answer-Choice Leverage}: All answers are odd $(5,6,7,8,9)$ but $6,8$ are even - suggests they're not actually possible, hinting at parity argument
    \item \textbf{Make It Easier}: Ignore the sum constraint initially, focus only on consecutive constraint
\end{itemize}

\subsection*{D. Argument Construction Strategy}
\begin{itemize}[leftmargin=*]
    \item \textbf{Direct Proof}: Use checkerboard coloring to show odd numbers must occupy specific positions, then apply sum constraint
    \item \textbf{Proof Structure}:
    \begin{enumerate}
        \item Establish that consecutive numbers alternate between "black" and "white" squares
        \item Count positions: 5 of one color, 4 of the other
        \item Since there are 5 odd numbers and 4 even numbers, odd numbers occupy the majority color
        \item Determine which positions are which color
        \item Apply the corner sum constraint to find the center value
    \end{enumerate}
\end{itemize}
\end{strategicbox}

\begin{tacticalbox}
\subsection*{A. Core Mathematical Tactics}
\begin{itemize}[leftmargin=*]
    \item \textbf{Symmetry}: The grid has rotational and reflectional symmetry which constrains valid arrangements
    \item \textbf{Parity Argument}: \textbf{PRIMARY TACTIC}
    \begin{itemize}
        \item Checkerboard coloring: Color squares alternately black/white
        \item Consecutive numbers must alternate colors (share edge, not diagonal)
        \item Therefore: Numbers $1,2,3,\ldots,9$ alternate: odd-even-odd-even-...
        \item Result: All odd numbers on one color, all even numbers on the other
    \end{itemize}
    \item \textbf{Invariant/Monovariant}: The alternating color property is invariant under the consecutive constraint
    \item \textbf{Graph Theory Perspective}: The arrangement forms a Hamiltonian path on the grid graph
\end{itemize}

\subsection*{B. Domain-Specific Tactics}
\begin{itemize}[leftmargin=*]
    \item \textbf{Combinatorics}: Count positions by type and apply pigeonhole principle
    \item \textbf{Number Theory}: Sum of all odd numbers $1+3+5+7+9=25$, sum of corners is $18$, so center is $25-18=7$
    \item \textbf{Position Analysis}: 
    \begin{itemize}
        \item 4 corner positions (2 neighbors each)
        \item 4 edge positions (3 neighbors each)  
        \item 1 center position (4 neighbors)
    \end{itemize}
\end{itemize}

\subsection*{C. Crossover Tactics}
\begin{itemize}[leftmargin=*]
    \item \textbf{Graph Theory}: Model as Hamiltonian path problem on $3\times3$ grid graph
    \item \textbf{Discrete Mathematics}: Bipartite graph coloring
    \item \textbf{Pattern Recognition}: Checkerboard pattern appears in many discrete problems
\end{itemize}
\end{tacticalbox}

\begin{conversionbox}
\subsection*{A. High-Probability Techniques (Recognition Index)}

\begin{itemize}[leftmargin=*]
    \item \textbf{Checkerboard/Parity Coloring} (90\% applicable)
    \begin{itemize}
        \item \textbf{Recognition}: "Consecutive numbers share an edge" $\to$ alternating pattern
        \item \textbf{Application}: Color grid black/white, observe numbers must alternate colors
        \item \textbf{Result}: 5 positions one color (odd numbers), 4 positions other color (even numbers)
    \end{itemize}
    
    \item \textbf{Complementary Counting/Sum Formula} (85\% applicable)
    \begin{itemize}
        \item \textbf{Recognition}: "Corner sum = 18" plus "all odd in corners/center"
        \item \textbf{Application}: Sum of all odd numbers = $1+3+5+7+9=25$
        \item \textbf{Formula}: Center = (Sum of all odds) - (Sum of corners) = $25-18=7$
    \end{itemize}
    
    \item \textbf{Extreme Value Analysis} (70\% applicable)
    \begin{itemize}
        \item \textbf{Recognition}: Numbers 1 and 9 must be at ends of the path
        \item \textbf{Application}: Since path has two endpoints, $1$ and $9$ must be at positions with fewer neighbors
        \item \textbf{Result}: $1$ and $9$ likely in corners (2 neighbors) not center (4 neighbors)
    \end{itemize}
\end{itemize}

\subsection*{B. Medium-Probability Techniques}
\begin{itemize}[leftmargin=*]
    \item \textbf{Constructive Verification} (60\% applicable)
    \begin{itemize}
        \item \textbf{Application}: After determining answer is 7, construct an actual valid arrangement
        \item \textbf{Example}: $\begin{array}{|c|c|c|}\hline 1 & 2 & 3 \\\hline 8 & 7 & 4 \\\hline 9 & 6 & 5 \\\hline\end{array}$
        \item \textbf{Verification}: Check corners: $1+3+9+5=18$ $\checkmark$, center is 7 $\checkmark$
    \end{itemize}
\end{itemize}

\subsection*{C. Technique Integration Strategy}
\begin{itemize}[leftmargin=*]
    \item \textbf{Primary Path}: Parity argument $\to$ Sum formula $\to$ Answer
    \item \textbf{Verification Path}: Constructive example
    \item \textbf{Time Allocation}: 
    \begin{itemize}
        \item Parity insight: 1 minute
        \item Sum calculation: 30 seconds
        \item Verification: 1 minute
    \end{itemize}
    \item \textbf{Confidence Boost}: Two independent methods confirm answer
\end{itemize}
\end{conversionbox}

\begin{verificationbox}
\subsection*{A. Verification Level Selection}

\textbf{Decision Tree Analysis:}
\begin{itemize}[leftmargin=*]
    \item Problem 12 → Mid-band → Default Level 3 (Standard)
    \item Parity argument is elegant and certain → Upgrade to Level 4 (Enhanced)
    \item Time available (3 min problem, spent 2 min) → Can afford Level 4
    \item \textbf{Selected Level}: Level 4 - Enhanced Verification
\end{itemize}

\subsection*{B. Verification Methods Applied}

\textbf{Level 4 - Enhanced Verification:}
\begin{enumerate}[leftmargin=*]
    \item \textbf{Arithmetic Verification}:
    \begin{itemize}
        \item Sum of odd numbers: $1+3+5+7+9=25$ $\checkmark$
        \item Corner constraint: $18$
        \item Center: $25-18=7$ $\checkmark$
    \end{itemize}
    
    \item \textbf{Parity Verification}:
    \begin{itemize}
        \item Checkerboard: 5 black, 4 white (or vice versa)
        \item 5 odd numbers, 4 even numbers
        \item Odd numbers on majority color $\checkmark$
        \item Center is majority color $\checkmark$
    \end{itemize}
    
    \item \textbf{Constructive Verification}:
    \begin{itemize}
        \item Build explicit example:
        \[
        \begin{array}{|c|c|c|}
        \hline
        1 & 2 & 3 \\
        \hline
        8 & 7 & 4 \\
        \hline
        9 & 6 & 5 \\
        \hline
        \end{array}
        \]
        \item Check consecutive adjacency: $1-2$, $2-3$, $3-4$, $4-5$, $5-6$, $6-7$, $7-8$, $8-9$ $\checkmark$
        \item Check corners: $1+3+9+5=18$ $\checkmark$
        \item Check center: $7$ $\checkmark$
    \end{itemize}
    
    \item \textbf{Answer Choice Elimination}:
    \begin{itemize}
        \item Can center be 5? Then corners sum to $25-5=20$ (need to be $18$) $\times$
        \item Can center be 6? Even number in center contradicts parity $\times$
        \item Can center be 7? Corners sum to $25-7=18$ $\checkmark$
        \item Can center be 8? Even number in center contradicts parity $\times$
        \item Can center be 9? Number 9 should be endpoint (corner), not center $\times$
    \end{itemize}
\end{enumerate}

\subsection*{C. Confidence Calibration}
\begin{itemize}[leftmargin=*]
    \item \textbf{Initial Confidence}: 85\% (after parity argument)
    \item \textbf{After Sum Verification}: 95\% (independent confirmation)
    \item \textbf{After Constructive Example}: 99\% (explicit construction works)
    \item \textbf{Final Confidence}: 99\% - extremely high confidence, answer is certainly (C) 7
\end{itemize}
\end{verificationbox}

\begin{roadmapbox}
\subsection*{Step-by-Step Solution Plan}

\textbf{Phase 1: Understanding (30 seconds)}
\begin{enumerate}[leftmargin=*]
    \item Read problem carefully, identify key constraints
    \item Recognize consecutive adjacency forms a path
    \item Note that all answer choices are in range 5-9
    \item \textbf{Checkpoint}: Clear understanding of constraints? Yes → Continue
\end{enumerate}

\textbf{Phase 2: Parity Insight (60 seconds)}
\begin{enumerate}[leftmargin=*]
    \item Visualize $3\times3$ grid with checkerboard coloring
    \item Observe: consecutive numbers must alternate colors (share edge)
    \item Count: 5 squares of one color, 4 of the other
    \item Deduce: 5 odd numbers $(1,3,5,7,9)$ on majority color, 4 even numbers $(2,4,6,8)$ on minority color
    \item Identify: Center is majority color → center must be odd
    \item \textbf{Checkpoint}: Eliminated even answers (6,8)? Yes → Continue
\end{enumerate}

\textbf{Phase 3: Position Analysis (45 seconds)}
\begin{enumerate}[leftmargin=*]
    \item Determine which positions are which color:
    \begin{itemize}
        \item If corners are black: 4 corners + 1 center = 5 black positions
        \item If corners are white: 4 edges + 1 center? No, that's only 5 positions but wrong configuration
    \end{itemize}
    \item Correct configuration: Corners and center are the same color (5 positions)
    \item Therefore: All odd numbers $\{1,3,5,7,9\}$ in corners and center
    \item \textbf{Checkpoint}: Confident odd numbers in corners/center? Yes → Continue
\end{enumerate}

\textbf{Phase 4: Sum Calculation (30 seconds)}
\begin{enumerate}[leftmargin=*]
    \item Sum of all odd numbers: $1+3+5+7+9=25$
    \item Given: Corner sum = $18$
    \item Calculate: Center = $25-18=7$
    \item \textbf{Checkpoint}: Arithmetic correct? Yes → Continue
\end{enumerate}

\textbf{Phase 5: Verification (45 seconds)}
\begin{enumerate}[leftmargin=*]
    \item Quick construction check: Can we build a valid grid with 7 in center?
    \item Try: $\begin{array}{ccc} 1 & 2 & 3 \\ 8 & 7 & 4 \\ 9 & 6 & 5 \end{array}$
    \item Verify: $1+3+9+5=18$ $\checkmark$, consecutive pairs share edges $\checkmark$
    \item \textbf{Checkpoint}: Construction valid? Yes → Select answer
\end{enumerate}

\subsection*{Alternative Approaches}

\textbf{Alternative 1: Brute Force Construction}
\begin{itemize}[leftmargin=*]
    \item Time: 4-5 minutes
    \item Method: Try to build valid arrangements systematically
    \item Risk: Time-consuming, may not finish
    \item Use case: If parity insight doesn't click
\end{itemize}

\textbf{Alternative 2: Answer Choice Testing}
\begin{itemize}[leftmargin=*]
    \item Time: 2-3 minutes
    \item Method: For each answer choice, check if valid arrangement exists
    \item Risk: Still time-consuming
    \item Use case: As verification after main solution
\end{itemize}

\subsection*{Common Pitfalls and Prevention}

\textbf{Pitfall 1: Misunderstanding "share an edge"}
\begin{itemize}[leftmargin=*]
    \item \textbf{Error}: Thinking diagonally adjacent squares satisfy the constraint
    \item \textbf{Prevention}: Explicitly note "share an edge" means horizontal/vertical only
    \item \textbf{Time cost if error}: 2-3 minutes to debug
\end{itemize}

\textbf{Pitfall 2: Incorrect parity analysis}
\begin{itemize}[leftmargin=*]
    \item \textbf{Error}: Not recognizing checkerboard pattern or miscounting colors
    \item \textbf{Prevention}: Draw the grid explicitly, count carefully
    \item \textbf{Time cost if error}: 1-2 minutes to correct
\end{itemize}

\textbf{Pitfall 3: Arithmetic error in sum}
\begin{itemize}[leftmargin=*]
    \item \textbf{Error}: Miscalculating $1+3+5+7+9$ or $25-18$
    \item \textbf{Prevention}: Double-check arithmetic, use verification
    \item \textbf{Time cost if error}: 30 seconds to fix
\end{itemize}

\textbf{Pitfall 4: Assuming answer without checking}
\begin{itemize}[leftmargin=*]
    \item \textbf{Error}: Jumping to answer without verifying it's achievable
    \item \textbf{Prevention}: Construct explicit example as verification
    \item \textbf{Time cost if error}: Could select wrong answer
\end{itemize}

\subsection*{Time Management}
\begin{itemize}[leftmargin=*]
    \item \textbf{Target Time}: 2.5-3 minutes
    \item \textbf{Actual Time Breakdown}:
    \begin{itemize}
        \item Understanding: 30 seconds
        \item Parity insight: 60 seconds
        \item Position analysis: 45 seconds
        \item Sum calculation: 30 seconds
        \item Verification: 45 seconds
        \item \textbf{Total}: 3 minutes 30 seconds
    \end{itemize}
    \item \textbf{Pivot Indicator}: If no insight after 2 minutes, try construction approach
\end{itemize}
\end{roadmapbox}

\begin{competitionbox}
\subsection*{A. Time Management Strategy}
\begin{itemize}[leftmargin=*]
    \item \textbf{Time Budget}: 3 minutes (Problem 12 of 25 in AMC12)
    \item \textbf{Time Allocation}:
    \begin{itemize}
        \item First pass (understanding): 30 sec
        \item Main solution: 2 min
        \item Verification: 30 sec
    \end{itemize}
    \item \textbf{Actual vs Expected}: May take 3-3.5 min with verification, acceptable for P12
\end{itemize}

\subsection*{B. Problem Prioritization}
\begin{itemize}[leftmargin=*]
    \item \textbf{Accessibility}: High - parity argument is clean once recognized
    \item \textbf{Domain Comfort}: Combinatorics/logic - suitable for most students
    \item \textbf{Recommendation}: Solve in first pass through mid-band problems (9-16)
\end{itemize}

\subsection*{C. Risk Management}
\begin{itemize}[leftmargin=*]
    \item \textbf{Confidence Level}: 99\% (after verification)
    \item \textbf{Error Recovery}: If parity approach fails, switch to construction/testing
    \item \textbf{Flag for Review}: No - very high confidence with multiple confirmations
\end{itemize}

\subsection*{D. Abandonment Criteria}
\begin{itemize}[leftmargin=*]
    \item \textbf{Soft Abandon}: If no progress after 4 minutes
    \item \textbf{Hard Abandon}: If construction becomes too messy or confusing
    \item \textbf{Strategy}: This problem unlikely to require abandonment given clean approach
\end{itemize}

\subsection*{E. Optimization Strategies}
\begin{itemize}[leftmargin=*]
    \item \textbf{Speed-up}: Recognize parity pattern immediately from "consecutive share edge"
    \item \textbf{Answer-choice leverage}: Note all choices are close to each other (5-9), suggests arithmetic constraint
    \item \textbf{Pattern recognition}: Checkerboard is common in AMC problems involving adjacency
\end{itemize}
\end{competitionbox}

\begin{keyinsightbox}
\textbf{Key Insight}: The consecutive adjacency constraint creates a checkerboard parity pattern. Since there are 5 odd numbers and 4 even numbers, and the grid has 5 positions of one color and 4 of the other, all odd numbers (including the center) must occupy the corners and center. The sum constraint then uniquely determines the center value as $25-18=7$.
\end{keyinsightbox}

\begin{warningbox}
\textbf{Warning}: Don't confuse "share an edge" with "adjacent" (which might include diagonals). The consecutive constraint is specifically horizontal/vertical adjacency only. Also, don't forget that the problem asks for the \textit{center} value, not a corner value.
\end{warningbox}

\begin{tipbox}
\textbf{Pro Tip}: Whenever you see "consecutive elements must be adjacent on a grid," immediately think about checkerboard coloring and parity arguments. This pattern appears frequently in AMC/AIME problems. Also, the elegant sum formula ($25-18=7$) provides a quick arithmetic check that can be done mentally.
\end{tipbox}

\section*{Final Answer}
\boxed{\textbf{(C)}\ 7}

\section*{Confidence Assessment}
\begin{itemize}[leftmargin=*]
    \item \textbf{Confidence Level}: 99\% - Parity argument is rigorous, sum calculation is straightforward, and constructive verification confirms the answer
    \item \textbf{Verification Applied}: Level 4 Enhanced - arithmetic verification, parity check, constructive example, and answer choice elimination
    \item \textbf{Flag for Review}: No - extremely high confidence with multiple independent confirmations
    \item \textbf{Time Spent}: Approximately 3 minutes 30 seconds (within acceptable range for Problem 12)
\end{itemize}

\section*{Solution Summary}

\textbf{Elegant Parity Solution:}

The key insight is recognizing that consecutive numbers must alternate between the two colors in a checkerboard pattern. Since the grid has 5 squares of one color and 4 of the other, and we have 5 odd numbers $(1,3,5,7,9)$ and 4 even numbers $(2,4,6,8)$, all odd numbers must occupy the five positions of one color.

When we color the grid like a checkerboard, the four corners and the center form five positions of the same color. Therefore, the odd numbers $\{1,3,5,7,9\}$ occupy exactly these five positions.

Since the corners sum to $18$ and all five odd numbers sum to $25$, the center must be:
$$25 - 18 = 7$$

This can be verified by constructing an explicit example:
$$\begin{array}{|c|c|c|}
\hline
1 & 2 & 3 \\
\hline
8 & 7 & 4 \\
\hline
9 & 6 & 5 \\
\hline
\end{array}$$

Corner sum: $1+3+9+5=18$ $\checkmark$, Center: $7$ $\checkmark$

\end{document}